% Hebrew abstract body — six labelled paragraphs mirroring the English
% structure. PRD V2 §7.2 Q5, FX-16. Operator reviews before D5 QA.

\noindent\textbf{רקע.}\quad
חיישני הדמיה בתחום הטרה-הרץ (THz) מפיקים תמונות בעלות רזולוציה נמוכה,
יחס אות-לרעש נמוך והעדר צבע. מסווגי \emph{ImageNet} שאומנו מראש,
המסתמכים על מרקמים עדינים ועל צבע רווי, מאבדים דיוק במהירות תחת תנאים
אלה, אך עוצמת הירידה והמנגנון העומד בבסיסה אינם מתועדים כמותית עבור
שלוש משפחות הארכיטקטורות המודרניות.

\noindent\textbf{מטרה.}\quad
להשוות את עמידותן של שלוש ארכיטקטורות-עומק מייצגות (רשת CNN שאריתית,
רשת CNN צפופה, וטרנספורמר ראייה פוֹביאַלי) תחת השפלה סינתטית מבוקרת
דמוית-THz, לפרק את ההפסד לתרומות לפי ציר, ולכמת כמה ממנו ניתן לשחזר על
ידי רגולריזציית-אימון.

\noindent\textbf{שיטות.}\quad
המערכה כוללת 402 תאי אימון של ResNet50, DenseNet121 ו-TransNeXt-tiny על
CIFAR-10 ו-MNIST, תחת עקומת השפלה בת חמש רמות (רוויה, הקטנת רזולוציה,
טשטוש, רעש גאוסי, מלח-ופלפל). שבעה שלבים מבודדים את האפקטים: A קו בסיס
נקי, B1 השפלה משולבת, C ציר יחיד, D שחזור באמצעות רגולריזציה, B2 ו-B2-nr
לפישוט פרוטוקול ה-THz, ו-C2 ציר יחיד תחת אותו פרוטוקול. בדיקת רבייה
על 48 חזרות עם זרעי-אקראיות נוספים שולחת חסם על שונות-בין-ריצות בכל אחד
מ-24 תאי הכותרת ברמה L3. כל צינור השפלה מעוגן לבדיקת PSNR / SSIM
דטרמיניסטית על 256 תמונות, כך שההשוואה בין השלבים מתבצעת על סקאלה
פיזיקלית.

\noindent\textbf{ממצאים.}\quad
תחת שלב B1 ב-CIFAR-10, TransNeXt-tiny שומר על המרווח הרחב ביותר מעל שתי
רשתות ה-CNN בקצה המתון של העקומה (כ-$10.1\,$\Mpp{} ברמה L1, כ-$10.4\,$\Mpp{}
ברמה L2 מול הממוצע של ה-CNN-ים), ומרווח זה מצטמצם לפחות מ-$1.5\,$\Mpp{}
החל מ-L3. ציר הרזולוציה דומיננטי בייחוס לפי ציר בקירוב פי ארבעה
(אפקט ממוצע L1 עד L5 על CIFAR-10: $-27.45\,$\Mpp{} עבור רזולוציה לעומת
$-6.58$ עבור טשטוש ו-$-5.13$ עבור מלח-ופלפל). רגולריזציה (T3 בשלב D)
משחזרת בקירוב $0.77\,$\Mpp{} מתוך קריסת הדיוק של $\sim 12.5\,$\Mpp{} בשלב
B1. בדיקת הרבייה מאשרת $\sigma < 1.5\,$\Mpp{} בכל תא כותרת ברמה L3.

\noindent\textbf{מסקנות.}\quad
קריסת הדיוק מתוארת היטב כצוואר-בקבוק של מידע ולא כהתאמת-יתר. הטרנספורמר
הפוֹביאַלי שומר על מרווח העמידות הגבוה ביותר בקצה המתון של עקומת ההשפלה;
שלוש הארכיטקטורות מתכנסות לרצועת $[19, 21]\,$\Mpp{} ברמה L5. פריסות THz
עתידיות צריכות להשקיע בשחזור רזולוציה לפני שלבי הסרת רעש או הסרת טשטוש.

\noindent\textbf{מילות מפתח.}\quad
סיווג עצמים, למידה עמוקה, רזולוציה נמוכה, הדמיית THz, עמידות, PSNR, SSIM,
כיול, טרנספורמר, רשת שאריתית, רשת צפופה.
